\documentclass[11pt]{article}
\usepackage[T1]{fontenc}
\usepackage{textcomp}
\usepackage[margin=0.75in]{geometry}
\usepackage{amsmath,amssymb,amsthm}
\usepackage{array,booktabs}
\usepackage{hyperref}
\setlength{\parindent}{0pt}
\newtheorem{theorem}{Theorem}
\theoremstyle{definition}
\newtheorem{definition}{Definition}
\title{Flow Proof Artifact: prop-32-tangent-chord-alternate-segment}
\author{Generated by \texttt{flow doc proof}}
\date{Flow Proof Book}
\begin{document}
\maketitle
If a straight line touches a circle and a chord is drawn from the point of contact, the angles with the tangent equal the angles in the alternate segments.\\[0.5em]
\textbf{Source.} euclid — Elements, Book III, Proposition 32\\[0.5em]
\bigskip
\noindent{\large \textbf{Derived fact 1.} \textit{If a straight line touches a circle and a chord is drawn from the point of contact, the angles with the tangent equal the angles in the alternate segments} \hfill the Euclidean plane \cdot Euclid Book III \cdot Proposition 32: the angle between tangent and chord equals the angle in the alternate segment}\par
\smallskip\noindent\textit{Coordinate: the Euclidean plane \cdot Euclid Book III \cdot Proposition 32: the angle between tangent and chord equals the angle in the alternate segment}\par
\smallskip\noindent\textit{Source: euclid — Elements, Book III, Proposition 32}\par
\smallskip\noindent\textit{Built on: proposition 19: the tangent is perpendicular to the diameter at the point of contact, for Euclid Book III on the Euclidean plane, proposition 21: angles in the same segment are equal, for Euclid Book III on the Euclidean plane}\par
\medskip
\noindent\textbf{Goal.} If a straight line touches a circle and a chord is drawn from the point of contact, the angles with the tangent equal the angles in the alternate segments\par
\noindent$\displaystyle \text{tangent chord angle equals alternate segment}$\par
\medskip
\noindent
\begin{tabular}{@{} >{\raggedright\arraybackslash}p{0.06\textwidth} >{\raggedright\arraybackslash}p{0.47\textwidth} >{\raggedleft\arraybackslash}p{0.41\textwidth} @{}}
\textbf{\#} & \textbf{Proof} & \textbf{Mathematics} \\
\midrule
\textbf{1.} & We prove that proposition 32: the angle between tangent and chord equals the angle in the alternate segment for Euclid Book III the Euclidean plane. & \\[0.45em]
\textbf{2.} & Let T = point\_of\_contact. & \\[0.45em]
\textbf{3.} & Let tangent = line\_through\_T. & \\[0.45em]
\textbf{4.} & Let chord = chord\_TA. & \\[0.45em]
\textbf{5.} & From step 2, step 3, and step 4, we can deduce that angle between tangent chord equals angle in alternate segment. & $\displaystyle \angle \text{between tangent chord} = \angle \text{in alternate segment}$ \\[0.45em]
\textbf{6.} & We invoke the derived fact governing Euclid Book III the Euclidean plane: proposition 19: the tangent is perpendicular to the diameter at the point of contact, for Euclid Book III on the Euclidean plane. & \\[0.45em]
\textbf{7.} & We invoke the derived fact governing Euclid Book III the Euclidean plane: proposition 21: angles in the same segment are equal, for Euclid Book III on the Euclidean plane. & \\[0.45em]
\textbf{8.} & From step 6 and step 7, this implies tangent chord angle equals alternate segment. Hence proven. & $\displaystyle \text{tangent chord angle equals alternate segment}$ \\[0.45em]
\bottomrule
\end{tabular}
\label{thm:Geometry-Euclid Book III-proposition_32_the_angle_between_tangent_and_chord_equals_the_angle_in_the_alternate_segment}
\par\medskip\hrule
\end{document}
